% !TEX root =./ECE343_Practice_main.tex
%
% L4 practice set (The Wave Equation): phasor form of the telegrapher's
% equations, the wave equation and its solution, and reading the solution.

\fullwidth{\textbf{Learning Objective 3.1: } I can describe the origin of the
transmission line equation, understand the method of developing a solution, and
interpret the solution.
\\[4pt]\normalfont\small Show your work. Write a \textbf{Documentation} line at the
end (\textbf{None} if you did not collaborate).}

\begin{questions}

%%%%%%%%%%%%%%%%%%%%%%%%%%%%%%%%%%%%%%%%%
\question \textbf{From the telegrapher's equations to the wave equation.} Start
from the telegrapher's equations,
\eq{
	-\partd{v(z,t)}{z} &= R^{\prime}\ i(z,t) + L^{\prime} \partd{i(z,t)}{t} \\
	-\partd{i(z,t)}{z} &= G^{\prime}\ v(z,t) + C^{\prime} \partd{v(z,t)}{t}
}
\begin{parts}

	\part With $v(z,t) = \text{Re}\left[\widetilde{V}(z) e^{j\omega t}\right]$,
	show that the phasor of $\partial v/\partial t$ is $j\omega\widetilde{V}(z)$.
	\begin{solution}[1.2in]
	\eq{
		\partd{v}{t} = \partd{}{t}\text{Re}\left[\widetilde{V}(z) e^{j\omega t}\right]
		= \text{Re}\left[\widetilde{V}(z) \partd{}{t} e^{j\omega t}\right]
		= \text{Re}\left[j\omega\widetilde{V}(z) e^{j\omega t}\right]
	}
	The derivative passes inside the real part because $t$ is real, and
	$\widetilde{V}(z)$ does not depend on $t$. The result is the real part of
	$j\omega\widetilde{V}(z)\,e^{j\omega t}$, so its phasor is $j\omega\widetilde{V}(z)$.
	\end{solution}

	\part Write both telegrapher's equations in phasor form.
	\begin{solution}[1in]
	\eq{
		-\dfrac{d\widetilde{V}(z)}{dz} &= \left(R^{\prime} + j\omega L^{\prime}\right)\widetilde{I}(z) \\
		-\dfrac{d\widetilde{I}(z)}{dz} &= \left(G^{\prime} + j\omega C^{\prime}\right)\widetilde{V}(z)
	}
	\end{solution}

	\part Uncouple the equations to derive the wave equation for the current,
	$\widetilde{I}(z)$. Show each step.
	\begin{solution}[1.6in]
	Differentiate the second equation with respect to $z$, then substitute the
	first for $d\widetilde{V}/dz$:
	\eq{
		-\dfrac{d^2\widetilde{I}(z)}{dz^2} &= \left(G^{\prime} + j\omega C^{\prime}\right)\dfrac{d\widetilde{V}(z)}{dz} \\
		&= -\left(G^{\prime} + j\omega C^{\prime}\right)\left(R^{\prime} + j\omega L^{\prime}\right)\widetilde{I}(z) \\
		\dfrac{d^2\widetilde{I}(z)}{dz^2} - \gamma^2\widetilde{I}(z) &= 0, \qquad
		\gamma^2 = \left(R^{\prime} + j\omega L^{\prime}\right)\left(G^{\prime} + j\omega C^{\prime}\right)
	}
	\end{solution}

	\part Substitute the trial solution $\widetilde{I}(z) = A e^{sz}$ into your
	result, find the characteristic equation and its roots, and write the general
	solution.
	\begin{solution}[1.3in]
	\eq{
		s^2 A e^{sz} - \gamma^2 A e^{sz} &= 0 \\
		s^2 - \gamma^2 &= 0 \quad\Longrightarrow\quad s = \pm\gamma \\
		\widetilde{I}(z) &= I_0^+ e^{-\gamma z} + I_0^- e^{+\gamma z}
	}
	\end{solution}

\end{parts}

\newpage
%%%%%%%%%%%%%%%%%%%%%%%%%%%%%%%%%%%%%%%%%
\question \textbf{The propagation constant of a lossy line.} A line has
$R^{\prime} = 0.5\ \ohms/\text{m}$, $L^{\prime} = 250\ \text{nH/m}$, $G^{\prime} = 0$, and
$C^{\prime} = 100\ \text{pF/m}$, and carries a $10\ \mhz$ signal.
\begin{parts}

	\begin{minipage}[t]{0.58\linewidth}
		\part Find the series impedance per meter, $R^{\prime} + j\omega L^{\prime}$, and the
		shunt admittance per meter, $G^{\prime} + j\omega C^{\prime}$.
		\begin{solution}[1in]
		\eq{
			\omega &= 2\pi\times10^{7} = 6.283\times10^{7}\ \text{rad/s} \\
			R^{\prime} + j\omega L^{\prime} &= 0.5 + j15.71\ \ohms/\text{m} \\
			G^{\prime} + j\omega C^{\prime} &= j6.283\times10^{-3}\ \text{S/m}
		}
		\end{solution}
	\end{minipage}\hfill%
	\begin{minipage}[t]{0.37\linewidth}
		\raggedleft \vspace{12pt}
		$Z^{\prime}$ = \ansbox[1.3in]{$0.5 + j15.71\ \ohms$/m}

		\vspace{4pt}
		$Y^{\prime}$ = \ansbox[1.3in]{$j6.28\times10^{-3}$ S/m}
	\end{minipage}

	\begin{minipage}[t]{0.58\linewidth}
		\part Find $\gamma = \alpha + j\beta$.
		\begin{solution}[1.6in]
		\eq{
			\gamma^2 &= (0.5 + j15.71)(j6.283\times10^{-3}) \\
			&= -0.0987 + j0.00314 \\
			&= 0.0987\angle 178.18^\circ \\
			\gamma &= \sqrt{0.0987}\angle 89.09^\circ = 0.3142\angle 89.09^\circ \\
			&= 0.0050 + j0.3142\ \text{m}^{-1}
		}
		We take the root with $\alpha \ge 0$ and $\beta > 0$.
		\end{solution}
	\end{minipage}\hfill%
	\begin{minipage}[t]{0.37\linewidth}
		\raggedleft \vspace{12pt}
		$\alpha$ = \ansbox[1.3in]{$0.0050$ Np/m}

		\vspace{4pt}
		$\beta$ = \ansbox[1.3in]{$0.314$ rad/m}
	\end{minipage}

	\begin{minipage}[t]{0.58\linewidth}
		\part Find the wavelength on the line.
		\begin{solution}[0.8in]
		\eq{ \lambda = \dfrac{2\pi}{\beta} = \dfrac{2\pi}{0.3142} = 20.0\ \m }
		\end{solution}
	\end{minipage}\hfill%
	\begin{minipage}[t]{0.37\linewidth}
		\raggedleft \vspace{12pt}
		$\lambda$ = \ansbox[1.3in]{$20.0\ \m$}
	\end{minipage}

	\part Compare $\beta$ with its value on the same line with $R^{\prime} = 0$. What does
	the comparison say about the effect of a small loss on the phase constant?
	\begin{solution}[1in]
	With $R^{\prime} = 0$, $\beta = \omega\sqrt{L^{\prime}C^{\prime}} = 6.283\times10^{7}\sqrt{2.5\times10^{-17}} = 0.3142$
	rad/m, the same to four figures. A small loss adds attenuation but leaves the
	phase constant, and so the wavelength, essentially unchanged.
	\end{solution}

\end{parts}

\newpage
%%%%%%%%%%%%%%%%%%%%%%%%%%%%%%%%%%%%%%%%%
\question \textbf{Reading a solution.} The voltage on a line is
\eq{
	v(z,t) = 5\,e^{-0.02z}\cosp{2\pi\times10^{8}\,t - 3z + 30^\circ}
	+ 2\,e^{0.02z}\cosp{2\pi\times10^{8}\,t + 3z}\ \text{V}
}
with $z$ in meters.
\begin{parts}

	\part Which term travels toward $+z$ and which toward $-z$? Explain how you
	can tell.
	\begin{solution}[1in]
	The first term travels toward $+z$: a point of constant phase,
	$\omega t - 3z = $ constant, moves to larger $z$ as $t$ grows. The second
	travels toward $-z$, because holding $\omega t + 3z$ constant requires $z$ to
	decrease as $t$ grows.
	\end{solution}

	\begin{minipage}[t]{0.58\linewidth}
		\part Find $\alpha$, $\beta$, and the wavelength.
		\begin{solution}[0.9in]
		\eq{
			\alpha = 0.02\ \text{Np/m}, \quad \beta = 3\ \text{rad/m}, \quad
			\lambda = \dfrac{2\pi}{3} = 2.09\ \m
		}
		\end{solution}
	\end{minipage}\hfill%
	\begin{minipage}[t]{0.37\linewidth}
		\raggedleft \vspace{12pt}
		$\alpha$ = \ansbox[1.3in]{$0.02$ Np/m}

		\vspace{4pt}
		$\lambda$ = \ansbox[1.3in]{$2.09\ \m$}
	\end{minipage}

	\begin{minipage}[t]{0.58\linewidth}
		\part Write the phasor constants $V_0^+$ and $V_0^-$.
		\begin{solution}[0.8in]
		\eq{ V_0^+ = 5\angle 30^\circ\ \text{V}, \qquad V_0^- = 2\angle 0^\circ\ \text{V} }
		\end{solution}
	\end{minipage}\hfill%
	\begin{minipage}[t]{0.37\linewidth}
		\raggedleft \vspace{12pt}
		$V_0^+$ = \ansbox[1.3in]{$5\angle 30^\circ$ V}

		\vspace{4pt}
		$V_0^-$ = \ansbox[1.3in]{$2\angle 0^\circ$ V}
	\end{minipage}

	\part The second term grows as $e^{0.02z}$. Explain why this does not mean the
	reflected wave gains energy.
	\begin{solution}[1in]
	The reflected wave travels toward $-z$. Following it in its own direction of
	travel, $z$ decreases, so $e^{0.02z}$ decreases: the wave decays as it
	travels, just as the incident wave does.
	\end{solution}

\end{parts}

\newpage
%%%%%%%%%%%%%%%%%%%%%%%%%%%%%%%%%%%%%%%%%
\question \textbf{The lossless line.} RG-58 has $L^{\prime} = 259\ \text{nH/m}$ and
$C^{\prime} = 96.8\ \text{pF/m}$.
\begin{parts}

	\part Set $R^{\prime} = G^{\prime} = 0$ and show that $\gamma$ is purely imaginary.
	\begin{solution}[1in]
	\eq{ \gamma = \sqrt{(j\omega L^{\prime})(j\omega C^{\prime})} = \sqrt{-\omega^2 L^{\prime}C^{\prime}} = j\omega\sqrt{L^{\prime}C^{\prime}} }
	so $\alpha = 0$ and $\beta = \omega\sqrt{L^{\prime}C^{\prime}}$.
	\end{solution}

	\begin{minipage}[t]{0.58\linewidth}
		\part Find $\beta$ and the wavelength at $50\ \mhz$.
		\begin{solution}[1.1in]
		\eq{
			\beta &= 2\pi\times5\times10^{7}\sqrt{259\times10^{-9}\times96.8\times10^{-12}} \\
			&= 1.57\ \text{rad/m} \\
			\lambda &= \dfrac{2\pi}{1.57} = 4.0\ \m
		}
		\end{solution}
	\end{minipage}\hfill%
	\begin{minipage}[t]{0.37\linewidth}
		\raggedleft \vspace{12pt}
		$\beta$ = \ansbox[1.3in]{$1.57$ rad/m}

		\vspace{4pt}
		$\lambda$ = \ansbox[1.3in]{$4.0\ \m$}
	\end{minipage}

	\part In the time domain, describe what the incident wave on a lossless line
	looks like to an observer at a fixed position $z$, and to an observer moving
	along with one of its crests.
	\begin{solution}[1in]
	At a fixed $z$, the voltage is a sinusoid in time with constant amplitude
	$\vert V_0^+\vert$. An observer moving with a crest sees a constant voltage,
	because the phase $\omega t - \beta z$ does not change, and with $\alpha = 0$
	the amplitude does not change either.
	\end{solution}

\end{parts}

\vspace{1.5em}
\noindent\textbf{Documentation:} \rule[-2pt]{0.6\linewidth}{0.4pt}

\end{questions}
